\documentclass[aspectratio=169,11pt]{beamer}
\usetheme{default}
\usecolortheme{default}

\setbeamertemplate{navigation symbols}{}
\setbeamertemplate{footline}[frame number]

\usepackage{amsmath}
\usepackage{booktabs}
\usepackage{graphicx}
\newcommand{\etal}{et al.}

\title{BZ Reaction--Diffusion: Transfer Functions, UQ, and a Search for the Right Training Task}
\author{Srivijayesh Venugopal}
\date{Cranfield IRP 2025-26 --- Supervisor Meeting, 11 August 2026}

\begin{document}

\begin{frame}
\titlepage
\end{frame}

% ---------------------------------------------------------------------
\begin{frame}{Meeting summary --- decision}

\textbf{Characterisation-first thesis.}

\begin{itemize}
    \item Claim: the structured light-sensitive Oregonator BZ medium is a \emph{quantifiable molecular-communication substrate}.
    \item Proof: verified solver + four transfer functions + Sobol/PCE sensitivity over $\varepsilon$, $f$, $\phi$, $D_u$.
    \item Training update: CMA-ES band-pass filter (208 evals) returned no improvement over uniform $\phi \approx 0.010$.
    \item Conclusion: the value is the \emph{datasheet and its uncertainty}, not a trained end-to-end circuit.
\end{itemize}

\vspace{0.3cm}
\textbf{Today:}
\begin{enumerate}
    \item Is this framing acceptable for submission?
    \item Should anisotropy stay in the body or move to appendix?
    \item If one quick task remains, chemical diode or none?
\end{enumerate}

\end{frame}

% ---------------------------------------------------------------------
\begin{frame}{Headline claim}

\textbf{Thesis claim:} characterise a structured BZ reaction--diffusion medium as a \emph{trainable analog substrate}, then quantify its native transfer functions and their parameter sensitivity \emph{before} attempting training.

\vspace{0.4cm}
\textbf{What we deliver today:}
\begin{enumerate}
    \item A verified Oregonator solver and four measured transfer functions.
    \item A first Sobol sensitivity / UQ pass over the kinetic parameters.
    \item An honest assessment of why the first training attempt was trivial.
    \item A proposal for a non-trivial, literature-backed training task we can still attempt.
\end{enumerate}

\vspace{0.3cm}
{\small The value is the datasheet, not a trained end-to-end circuit.}

\end{frame}

% ---------------------------------------------------------------------
\begin{frame}{Four measured transfer functions}

\begin{table}
\centering
\footnotesize
\begin{tabular}{p{3.2cm}p{4.5cm}p{5.5cm}}
\toprule
\textbf{Function} & \textbf{What it measures} & \textbf{Oregonator-A headline} \\
\midrule
Channel wire & Pulse speed and attenuation vs width & 6.22 cells/t.u., width-independent down to $W=1$ \\
Frequency response & Refractory bandwidth & 3.10 t.u. refractory; 0.282/t.u. max 1:1 rate \\
Collision logic & Native two-input gate & A AND (NOT B); true peak 0.708, false peaks $\sim$0.003 (rest-state floor), 232$\times$ separation \\
Anisotropic routing & Directional steering & $\sqrt{r}$ law verified; patch delay 644--836 steps \\
\bottomrule
\end{tabular}
\end{table}

\vspace{0.2cm}
{\small All four are native properties of the Oregonator excitable medium at the working point. False logic states sit at the rest-state noise floor ($\sim$0.003), not at exactly zero.}

\end{frame}

% ---------------------------------------------------------------------
\begin{frame}{What the transfer functions tell us}

\begin{itemize}
    \item \textbf{Channel:} walls are just wiring --- width carries almost no computing leverage.
    \item \textbf{Frequency:} refractoriness is a built-in low-pass filter; the medium has a hard clock ceiling.
    \item \textbf{Logic:} collisions annihilate; the native Boolean primitive is \textbf{inhibition}, not AND/OR/XOR directly.
    \item \textbf{Anisotropy:} the only graded, non-blocking routing knob, but experimentally harder to realise than $\phi(x,y)$.
\end{itemize}

\vspace{0.3cm}
\textbf{Bottom line:} the substrate's ``instruction set'' is small and constrained. Training can only reshuffle these primitives.

\end{frame}

% ---------------------------------------------------------------------
\begin{frame}{Uncertainty quantification: first Sobol pass}

Four parameters perturbed $\pm 10\%$ around the working point; 384 forward solves.

\vspace{0.3cm}
\textbf{Pulse speed:} dominated by $\varepsilon$ ($S_1 \approx 0.59$), then $D_u$ ($\approx 0.32$), $\phi$ ($\approx 0.11$), $f$ ($\approx 0.08$).

\vspace{0.2cm}
\textbf{Probe peak:} dominated by $f$ ($\approx 0.65$), then $\varepsilon$ ($\approx 0.37$), with $\phi$ and $D_u$ almost inactive.

\vspace{0.3cm}
\textbf{Take-away:} speed and amplitude have different dominant parameters; a calibration or training loop must target the right QoI. This is the UQ step that should precede any trained BZ experiment.

\end{frame}

% ---------------------------------------------------------------------
\begin{frame}{Training is now a feasibility/negative-result note}

\textbf{What we tried:} CMA-ES band-pass filter, 8 control segments, 208 evaluations.
\begin{itemize}
    \item Best loss $\approx 1.0000$; optimised control stayed at $\phi \approx 0.010$ everywhere.
    \item No better solution than the uniform channel found within this budget.
\end{itemize}

\vspace{0.3cm}
\textbf{Why this is informative, not a failure:}
\begin{itemize}
    \item The transfer functions already showed refractoriness is a hard ceiling.
    \item If the uniform channel is already near-optimal for rate shaping, the optimizer correctly returns it.
    \item It bounds how much $\phi(x,y)$ can reshape a single-channel pulse train.
\end{itemize}

\end{frame}

% ---------------------------------------------------------------------
\begin{frame}{If one task remains: trainable chemical diode}

The diode is the best remaining hedge if we want a trained-device chapter:
\begin{itemize}
    \item Input from either end of a channel; output at the opposite end.
    \item Target: high transmission one way, low transmission the other.
    \item Lever: asymmetric $\phi(x)$ profile along the channel.
    \item Literature precedent: Agladze \etal{} 1996; G\'{o}recka \etal{} 2009.
\end{itemize}

\vspace{0.3cm}
\textbf{Why it is sharper than the band-pass:}
\begin{itemize}
    \item Requires \emph{asymmetric dynamical response}, not just a static barrier.
    \item Single channel, 1-D $\phi$ profile --- cheap enough for a 1--2 day CMA-ES run.
    \item Success or failure both teach us something about the substrate's trainable degrees of freedom.
\end{itemize}

\end{frame}

% ---------------------------------------------------------------------
\begin{frame}{Why this fits your research}

Your work treats molecular communication and diffusion/reaction as information channels (Farsad \etal{} COMST 2016; Lin \etal{} IEEE Trans. Commun. 2023); Parrilla-Gutierrez \etal{} (Nat. Commun. 2020) showed a real programmable chemical computer.

\vspace{0.3cm}
This thesis takes the next step:
\begin{itemize}
    \item The channel is not only carrying messages --- its geometry and property fields \emph{are} the program.
    \item The measured transfer functions are the instruction set of a spatial, analog, chemical computer.
    \item The limitations (refractory clock, OR background at junctions, bounded $\phi$ window) are exactly the constraints a channel designer must respect.
\end{itemize}

\end{frame}

% ---------------------------------------------------------------------
\begin{frame}{Questions for you}

\begin{enumerate}
    \item Does the ``characterisation + UQ + one trained device'' framing read as a coherent MSc contribution?
    \item Which training task should I attempt in the remaining days, if at all: chemical diode or something else?
\end{enumerate}

\end{frame}

% ---------------------------------------------------------------------
\begin{frame}{Timeline to submission}

\begin{itemize}
    \item Thesis draft: clean compile; Sobol/UQ section in place.
    \item Training update: band-pass CMA-ES completed; no improvement over uniform $\phi$.
    \item Pending decisions:
    \begin{itemize}
        \item Accept characterisation + UQ + negative training result as submission framing?
        \item Demote anisotropy to appendix / future work?
        \item Authorise one 1--2 day chemical-diode run, or stop training now?
    \end{itemize}
\end{itemize}
\end{frame}

\end{document}
